\section{Project Plan}
\subsection{Process and Phase Gates}

The project follows the waterfall model, with test planning prepared in the
spirit of the V-model: the PA4 test plan is derived from the PA2 requirements
and the PA3 design. Each phase ends with a phase gate, which is a PA submission.
A submitted PA becomes a baseline; any later change is recorded, its impact is
assessed, and it is listed in the next PA. Functional coding is planned to
start after the PA3 gate.

\begin{table}[H]
\centering
\caption{Phase gates and exit criteria.}
\label{tab:gates}
\begin{tabularx}{\textwidth}{@{}L{1.2cm}L{3.0cm}L{2.6cm}Y@{}}
\toprule
\textbf{Gate} & \textbf{Phase} & \textbf{Planned due} & \textbf{Exit criteria}\\
\midrule
PA1 & Initiation and planning & End of Week 3 &
  Vision, scope, and the approved WBS baseline, estimates, schedule, and risks
  are consistent. GitHub Projects lists all current work packages, owners, and
  deadlines. Before PA1 submission, configure and verify main-branch protection
  and at least one required PR approval, grant and verify instructor/TA repository
  access, and confirm TA participation in Discord.\\
PA2 & Requirements & End of Week 5 &
  Functional and non-functional requirements, business rules and acceptance
  criteria are clear and prioritised; use cases and prototype are complete;
  the traceability matrix and requirements review are finished.\\
PA3 & Design & End of Week 7 &
  C4 diagrams, UML models, ERD and UI design are consistent with each other
  and with PA2; at least three ADRs are recorded; the design covers all
  Must-have requirements; major review findings are closed.\\
PA4 & Test planning & End of Week 8 &
  The test plan applies equivalence partitioning, boundary value analysis and
  decision tables; test cases trace to every Must-have requirement;
  the test environment and pass/fail criteria are defined.\\
PA5 & Implementation, testing and deployment & End of Week 10 &
  All Must-have features work and are deployed; automated tests cover at least
  two use cases with at least two scenarios per use case; test and coverage
  reports, CI, code-review evidence, working-software video, and user guide are
  available.\\
Final & Presentation and demo & Week 11 &
  A rehearsal of the demo succeeds; the slides are complete; every member
  speaks for at least 2 minutes and can explain their own work.\\
\bottomrule
\end{tabularx}
\end{table}

Week numbers are the planning baseline; official submission dates and
times are those published on Moodle.

\textbf{Change control.} Anyone may propose a change by opening a
change-request (CR) issue that states the reason, the affected requirements,
the estimated effort and the impact on the schedule. The business analyst
assesses the business impact, a developer assesses the technical impact, and
the team lead consolidates the results. The whole team decides; changes to the
committed scope are discussed with the TA. After approval, the documents, task
board and traceability matrix are updated, and the change is listed in the
next PA.

\subsection{Work Breakdown Structure}

The WBS has three levels. Level~1 consists of the project phases, level~2 of
deliverables or groups of deliverables, and level~3 of work packages. For
example: 5~Implementation, testing and deployment $\rightarrow$
5.1~Back end core and infrastructure $\rightarrow$
5.1.2~Booking and double-booking prevention. Each work package has exactly one
owner, an effort estimate in person-hours and a target week, and is tracked as
one issue on the task board. WBS identifiers are independent of the section
numbers of this document.

\begin{table}[H]
\centering
\caption{WBS levels 1 and 2.}
\label{tab:wbs-levels}
\begin{tabularx}{\textwidth}{@{}L{4.6cm}Y@{}}
\toprule
\textbf{Level 1 (phase)} & \textbf{Level 2 (deliverable groups)}\\
\midrule
1 Planning (PA1) &
  1.1 Vision; 1.2 Solution and work breakdown; 1.3 Estimation and schedule;
  1.4 Risk, tools and submission\\
2 Requirements (PA2) &
  2.1 Requirements specification; 2.2 Use cases and prototype;
  2.3 Review and baseline\\
3 Design (PA3) &
  3.1 Architecture and data; 3.2 Object-oriented and UI design;
  3.3 Review and packaging\\
4 Test planning (PA4) &
  4.1 Test design; 4.2 Review and environment\\
5 Implementation, testing and deployment (PA5) &
  5.1 Back end core and infrastructure; 5.2 Front end;
  5.3 Admin API and integration; 5.4 System testing and reporting\\
6 Final presentation &
  6.1 Presentation\\
7 Continuous management &
  7.1 Management throughout the project\\
\bottomrule
\end{tabularx}
\end{table}

Effort values are expert-judgement estimates in person-hours; further
estimation details are provided in Section~5.3. The effort of each phase
already includes its reviews and its business and progress meetings.
Continuous management (group 7) covers only the consolidation of weekly
reports and the maintenance of the task board and repository, so no effort is
counted twice.

\begin{longtable}{@{}L{1.4cm}L{7.4cm}L{1.6cm}rL{1.7cm}@{}}
\caption{Work packages for PA1 and PA2.}\label{tab:wbs-pa12}\\
\toprule
\textbf{WBS} & \textbf{Work package} & \textbf{Owner} & \textbf{Hours} & \textbf{Week}\\
\midrule
\endfirsthead
\toprule
\textbf{WBS} & \textbf{Work package} & \textbf{Owner} & \textbf{Hours} & \textbf{Week}\\
\midrule
\endhead
\bottomrule
\endfoot
\multicolumn{5}{@{}l}{\textbf{1\quad Planning (PA1)}}\\
\addlinespace[2pt]
\multicolumn{5}{@{}l}{\textit{1.1 Vision}}\\
1.1.1 & Vision and scope & Long & 4 & W3\\
1.1.2 & Stakeholders and product positioning & Long & 2 & W3\\
\multicolumn{5}{@{}l}{\textit{1.2 Solution and work breakdown}}\\
1.2.1 & Context diagram, technology stack and environment & Th\^ong & 4 & W3\\
1.2.2 & Phase gates and WBS & Th\^ong & 3 & W3\\
\multicolumn{5}{@{}l}{\textit{1.3 Estimation and schedule}}\\
1.3.1 & Estimation, Gantt schedule, roles and cost & Ng\^an & 5 & W3\\
\multicolumn{5}{@{}l}{\textit{1.4 Risk, tools and submission}}\\
1.4.1 & Risk register & Ph\'uc & 2 & W3\\
1.4.2 & Repository, task board and communication setup & Ph\'uc & 2 & W3\\
1.4.3 & Assembly, cross-review and submission of PA1 & Long & 2 & W3\\
\midrule
\multicolumn{3}{@{}l}{\textbf{Subtotal PA1}} & \textbf{24} & \\
\midrule
\multicolumn{5}{@{}l}{\textbf{2\quad Requirements (PA2)}}\\
\addlinespace[2pt]
\multicolumn{5}{@{}l}{\textit{2.1 Requirements specification}}\\
2.1.1 & Functional requirements and business rules & Th\^ong & 7 & W4--5\\
2.1.2 & Non-functional requirements and acceptance criteria & Long & 4 & W4--5\\
\multicolumn{5}{@{}l}{\textit{2.2 Use cases and prototype}}\\
2.2.1 & Use cases and exception flows & Th\^ong & 5 & W4--5\\
2.2.2 & Prototype of the main flows & Ng\^an & 8 & W4--5\\
\multicolumn{5}{@{}l}{\textit{2.3 Review and baseline}}\\
2.3.1 & Traceability matrix and SRS review & Ph\'uc & 6 & W5\\
2.3.2 & Assembly and baseline of PA2 & Long & 6 & W5\\
2.3.3 & Prototype review & Ng\^an & 2 & W5\\
2.3.4 & Data feasibility check & Ph\'uc & 2 & W4\\
\midrule
\multicolumn{3}{@{}l}{\textbf{Subtotal PA2}} & \textbf{40} & \\
\end{longtable}

\begin{longtable}{@{}L{1.4cm}L{7.4cm}L{1.6cm}rL{1.7cm}@{}}
\caption{Work packages for PA3 and PA4.}\label{tab:wbs-pa34}\\
\toprule
\textbf{WBS} & \textbf{Work package} & \textbf{Owner} & \textbf{Hours} & \textbf{Week}\\
\midrule
\endfirsthead
\toprule
\textbf{WBS} & \textbf{Work package} & \textbf{Owner} & \textbf{Hours} & \textbf{Week}\\
\midrule
\endhead
\bottomrule
\endfoot
\multicolumn{5}{@{}l}{\textbf{3\quad Design (PA3)}}\\
\addlinespace[2pt]
\multicolumn{5}{@{}l}{\textit{3.1 Architecture and data}}\\
3.1.1 & C4 architecture and ADRs & Ph\'uc & 6 & W6\\
3.1.2 & ERD and table specifications & Ph\'uc & 6 & W6--7\\
\multicolumn{5}{@{}l}{\textit{3.2 Object-oriented and UI design}}\\
3.2.1 & UML models and class specifications & Th\^ong & 8 & W6--7\\
3.2.2 & UI design and screen flow & Ng\^an & 8 & W6--7\\
\multicolumn{5}{@{}l}{\textit{3.3 Review and packaging}}\\
3.3.1 & Traceability update and design review & Long & 6 & W7\\
3.3.2 & Assembly of PA3 & Ng\^an & 2 & W7\\
3.3.3 & Final check and submission of PA3 & Long & 2 & W7\\
3.3.4 & Testability review & Th\^ong & 2 & W7\\
\midrule
\multicolumn{3}{@{}l}{\textbf{Subtotal PA3}} & \textbf{40} & \\
\midrule
\multicolumn{5}{@{}l}{\textbf{4\quad Test planning (PA4)}}\\
\addlinespace[2pt]
\multicolumn{5}{@{}l}{\textit{4.1 Test design}}\\
4.1.1 & Test plan, equivalence partitioning, boundary values, decision tables & Th\^ong & 4 & W8\\
4.1.2 & UI test cases & Ng\^an & 2 & W8\\
\multicolumn{5}{@{}l}{\textit{4.2 Review and environment}}\\
4.2.1 & Review and submission of PA4 & Long & 1 & W8\\
4.2.2 & Test environment and test traceability & Long & 1 & W8\\
\midrule
\multicolumn{3}{@{}l}{\textbf{Subtotal PA4}} & \textbf{8} & \\
\end{longtable}

The 8 hours of PA4 reuse the specifications and acceptance criteria produced
in PA2; the execution of system tests is counted under group 5.4.

\begin{longtable}{@{}L{1.4cm}L{7.4cm}L{1.6cm}rL{1.7cm}@{}}
\caption{Work packages for PA5, the final presentation and continuous
management.}\label{tab:wbs-pa5}\\
\toprule
\textbf{WBS} & \textbf{Work package} & \textbf{Owner} & \textbf{Hours} & \textbf{Week}\\
\midrule
\endfirsthead
\toprule
\textbf{WBS} & \textbf{Work package} & \textbf{Owner} & \textbf{Hours} & \textbf{Week}\\
\midrule
\endhead
\bottomrule
\endfoot
\multicolumn{5}{@{}l}{\textbf{5\quad Implementation, testing and deployment (PA5)}}\\
\addlinespace[2pt]
\multicolumn{5}{@{}l}{\textit{5.1 Back end core and infrastructure}}\\
5.1.1 & Database and authentication (F01) & Ph\'uc & 9 & W8\\
5.1.2 & Booking and double-booking prevention (F03, F04); review the PA3 booking API contract and acceptance criteria in W8, then implement booking and conflict prevention in W9 & Ph\'uc & 10 & W8--9\\
5.1.3 & Back end unit tests and CI; initial CI setup in W8, with unit tests and coverage checks completed in W10 & Ph\'uc & 5 & W8--10\\
5.1.4 & Deployment and smoke test & Ph\'uc & 3 & W10\\
\multicolumn{5}{@{}l}{\textit{5.2 Front end}}\\
5.2.1 & Room list and calendar (F02) & Ng\^an & 5 & W8\\
5.2.2 & Submit, view and cancel requests (F03) & Ng\^an & 7 & W8--9\\
5.2.3 & Administrator screens (F04, F05); prioritize approval UI for W9 integration and complete room/maintenance management in W10 & Ng\^an & 8 & W9--10\\
5.2.4 & Front end unit tests & Ng\^an & 3 & W9--10\\
\multicolumn{5}{@{}l}{\textit{5.3 Admin API and integration}}\\
5.3.1 & Approval and room-management API & Long & 6 & W8--9\\
5.3.2 & Integrate ready student-booking and admin-approval flows in W9; integrate room/maintenance management after its UI and API are complete in W10 & Long & 5 & W9--10\\
5.3.3 & Defect fixing after system testing & Long & 4 & W10\\
5.3.4 & Prepare document structure and a guide draft from available flows in W9; update against the tested deployed build and record the video in W10 & Long & 5 & W9--10\\
\multicolumn{5}{@{}l}{\textit{5.4 System testing and reporting}}\\
5.4.1 & Test data and test accounts & Th\^ong & 4 & W8\\
5.4.2 & Prepare and execute automated tests for at least two use cases, with at least two scenarios per use case, after the corresponding workflows are integrated & Th\^ong & 6 & W8--10\\
5.4.3 & Test integrated workflows in W9; complete system testing and retesting in W10 & Th\^ong & 9 & W9--10\\
5.4.4 & Prepare the report structure and collect available results in W9; finalize after W10 test results & Th\^ong & 5 & W9--10\\
\midrule
\multicolumn{3}{@{}l}{\textbf{Subtotal PA5}} & \textbf{94} & \\
\midrule
\multicolumn{5}{@{}l}{\textbf{6\quad Final presentation}}\\
\addlinespace[2pt]
\multicolumn{5}{@{}l}{\textit{6.1 Presentation}}\\
6.1.1 & Presentation script and overview slides & Long & 3 & W11\\
6.1.2 & Business demo and Q\&A & Th\^ong & 3 & W11\\
6.1.3 & UI slides and rehearsal & Ng\^an & 3 & W11\\
6.1.4 & Demo check and backup & Ph\'uc & 3 & W11\\
\midrule
\multicolumn{3}{@{}l}{\textbf{Subtotal Final}} & \textbf{12} & \\
\midrule
\multicolumn{5}{@{}l}{\textbf{7\quad Continuous management}}\\
\addlinespace[2pt]
\multicolumn{5}{@{}l}{\textit{7.1 Management throughout the project}}\\
7.1.1 & Coordination and weekly reports & Long & 9 & W2--10\\
7.1.2 & Board updates for business work & Th\^ong & 3 & W2--10\\
7.1.3 & Board updates for UI work & Ng\^an & 3 & W2--10\\
7.1.4 & Repository and technical progress updates & Ph\'uc & 3 & W2--10\\
\midrule
\multicolumn{3}{@{}l}{\textbf{Subtotal continuous management}} & \textbf{18} & \\
\end{longtable}

\textbf{Planned scenarios for WBS 5.4.2.} This is proposed test scope, not a
report of completed tests: submit booking request (valid request; request
outside opening hours) and approve booking request (available room/time;
conflicting booking, including the concurrent-approval case).

\begin{table}[H]
\centering
\caption{Effort summary by phase (baseline, before contingency).}
\label{tab:effort}
\begin{tabular}{@{}lr@{}}
\toprule
\textbf{Phase} & \textbf{Effort (person-hours)}\\
\midrule
1 Planning (PA1) & 24\\
2 Requirements (PA2) & 40\\
3 Design (PA3) & 40\\
4 Test planning (PA4) & 8\\
5 Implementation, testing and deployment (PA5) & 94\\
6 Final presentation & 12\\
7 Continuous management & 18\\
\midrule
\textbf{Total baseline effort} & \textbf{236}\\
\bottomrule
\end{tabular}
\end{table}

\textbf{Definition of done.} A work package is complete when:
\begin{itemize}[leftmargin=*,itemsep=2pt]
  \item the deliverable or code change described in its issue exists and its
        acceptance criteria are met;
  \item the designated reviewer has checked it, the documents or code are
        updated in the repository through a pull request, and no blocking
        comment remains;
  \item the owner has recorded the actual effort, the issue status, the link
        to the pull request or document, and any remaining issues on the task
        board.
\end{itemize}

\textbf{Early technical checkpoints.}
\begin{itemize}[leftmargin=*,itemsep=2pt]
  \item F03 and F04 share the same status model but are used by different
        roles. The sequence diagrams must show the permission check and the
        conflict check at approval time.
  \item The ERD and the transaction design must support rooms, bookings, user
        roles, maintenance periods, the fixed booking rules and the status
        history.
  \item Test design must cover time boundaries, capacity limits, bookings
        outside opening hours, cancellations close to the start time and
        concurrent approvals.
  \item Candidate ADRs for PA3 are the choice of PostgreSQL, the use of
        Firebase for authentication, and the mechanism that prevents two
        Approved bookings from overlapping. PA1 does not replace these ADRs.
\end{itemize}


\subsection{Estimation}

The team used bottom-up estimation and expert judgement. Package owners
estimated the work packages, and the team reviewed and approved the work
packages, owners, and estimates as the PA1 planning baseline. Effort is measured
in person-hours, not calendar duration; estimates include only the work and
review/test activities specified in each package, avoiding double-counting.
The approved estimated-effort baseline is 236 person-hours: PA1 24, PA2 40,
PA3 40, PA4 8, PA5 94, Final 12, and continuous management 18. These are
estimates, not actual hours worked. Changes after this baseline follow change
control.

\textbf{Planned work-package allocation for W8--W10.} The table gives the
planned weekly distribution in person-hours. Allocations for PA4 and PA5
packages sum to each package's full WBS effort. Rows marked with an asterisk
show only the W8--W10 share of the existing continuous-management packages.

\begingroup
\small
\setlength{\tabcolsep}{4pt}
\renewcommand{\arraystretch}{1.08}
\setlength{\LTleft}{\fill}
\setlength{\LTright}{\fill}
\begin{longtable}{@{}L{1.35cm}L{2.15cm}>{\centering\arraybackslash}p{1.15cm}>{\centering\arraybackslash}p{1.15cm}>{\centering\arraybackslash}p{1.15cm}@{}}
\caption{Planned W8--W10 work-package allocation (person-hours).}\label{tab:w810-wbs-allocation}\\
\toprule
\textbf{WBS ID} & \textbf{Owner} & \textbf{W8} & \textbf{W9} & \textbf{W10} \\
\midrule
\endfirsthead
\toprule
\textbf{WBS ID} & \textbf{Owner} & \textbf{W8} & \textbf{W9} & \textbf{W10} \\
\midrule
\endhead
\midrule
\endfoot
\bottomrule
\endlastfoot
4.1.1 & Th\^ong & 4 & 0 & 0 \\
4.1.2 & Ng\^an & 2 & 0 & 0 \\
4.2.1 & Long & 1 & 0 & 0 \\
4.2.2 & Long & 1 & 0 & 0 \\
5.1.1 & Ph\'uc & 9 & 0 & 0 \\
5.1.2 & Ph\'uc & 0.25 & 9.75 & 0 \\
5.1.3 & Ph\'uc & 0.5 & 0 & 4.5 \\
5.1.4 & Ph\'uc & 0 & 0 & 3 \\
5.2.1 & Ng\^an & 5 & 0 & 0 \\
5.2.2 & Ng\^an & 2.5 & 4.5 & 0 \\
5.2.3 & Ng\^an & 0 & 4.5 & 3.5 \\
5.2.4 & Ng\^an & 0 & 0.5 & 2.5 \\
5.3.1 & Long & 3 & 3 & 0 \\
5.3.2 & Long & 0 & 3 & 2 \\
5.3.3 & Long & 0 & 0 & 4 \\
5.3.4 & Long & 0 & 2 & 3 \\
5.4.1 & Th\^ong & 4 & 0 & 0 \\
5.4.2 & Th\^ong & 1 & 4 & 1 \\
5.4.3 & Th\^ong & 0 & 3 & 6 \\
5.4.4 & Th\^ong & 0 & 2.5 & 2.5 \\
\midrule
7.1.1* & Long & 1 & 1 & 1 \\
7.1.2* & Th\^ong & 0.25 & 0.25 & 0.5 \\
7.1.3* & Ng\^an & 0.25 & 0.25 & 0.5 \\
7.1.4* & Ph\'uc & 0.25 & 0.25 & 0.5 \\
\midrule
\multicolumn{2}{@{}l}{\textbf{W8--W10 total}} & \textbf{35} & \textbf{38.5} & \textbf{34.5} \\
\end{longtable}
\endgroup

The asterisked entries are the W8--W10 portions only; the full WBS effort
remains 9 hours for 7.1.1 and 3 hours for each of 7.1.2--7.1.4. These
allocations are planning figures, not actual effort or evidence of completed
work. The W8--W10 workload totals 108 hours: 94 for PA5, 8 for PA4, and 6
from existing continuous-management effort.

The usual planning allowance is about 7 hours per person per week. The team
has agreed to make 10 hours per person per week available in W8--W10, for a
confirmed capacity of $4 \times 10 \times 3 = 120$ person-hours. Long has
confirmed enough availability for his assigned W3 and W7 work, including
continuous-management tasks; the usual 7-hour allowance is not a hard limit for
his workload in those weeks.

\begin{table}[H]
\centering
\caption{W8--W10 planned effort and capacity by member (person-hours).}
\label{tab:w810-capacity}
\small
\setlength{\tabcolsep}{4pt}
\renewcommand{\arraystretch}{1.15}
\begin{tabular}{@{}lrrrrr@{}}
\toprule
\textbf{Member} & \textbf{PA5} & \textbf{PA4} & \textbf{\shortstack{Management\\W8--W10}} & \textbf{Total} & \textbf{Capacity} \\
\midrule
Long & 20 & 2 & 3 & 25 & 30 \\
Th\^ong & 24 & 4 & 1 & 29 & 30 \\
Ng\^an & 23 & 2 & 1 & 26 & 30 \\
Ph\'uc & 27 & 0 & 1 & 28 & 30 \\
\midrule
\textbf{Total} & \textbf{94} & \textbf{8} & \textbf{6} & \textbf{108} & \textbf{120} \\
\bottomrule
\end{tabular}
\end{table}

\begin{table}[H]
\centering
\caption{W8--W10 planned effort by member and week (person-hours).}
\label{tab:w810-weekly-load}
\small
\setlength{\tabcolsep}{6pt}
\renewcommand{\arraystretch}{1.15}
\begin{tabular}{@{}lrrrr@{}}
\toprule
\textbf{Member} & \textbf{W8} & \textbf{W9} & \textbf{W10} & \textbf{Total} \\
\midrule
Long & 6 & 9 & 10 & 25 \\
Th\^ong & 9.25 & 9.75 & 10 & 29 \\
Ng\^an & 9.75 & 9.75 & 6.5 & 26 \\
Ph\'uc & 10 & 10 & 8 & 28 \\
\midrule
\textbf{Total} & \textbf{35} & \textbf{38.5} & \textbf{34.5} & \textbf{108} \\
\bottomrule
\end{tabular}
\end{table}

The 12-hour difference between the agreed 120-hour capacity and the
108-hour workload is capacity headroom, not additional effort; the project
baseline remains 236 hours. Headroom is 5 hours for Long, 1 for Th\^ong, 4 for
Ng\^an, and 2 for Ph\'uc. W9 and W10 both include members at or near 10 hours,
so the PM should track per-member progress and handoffs rather than rely only on
group capacity.

\subsection{Schedule}

The schedule is a week-level planning baseline; official submission dates
and times are those published on Moodle. Tasks may overlap where dependencies
allow. The W8--W10 package allocations in Table~\ref{tab:w810-wbs-allocation}
provide the detailed weekly plan; the Gantt below summarizes those spans.

\begingroup
\small
\setlength{\tabcolsep}{2.5pt}
\renewcommand{\arraystretch}{1.08}
\setlength{\LTleft}{\fill}
\setlength{\LTright}{\fill}
\begin{longtable}{@{}>{\raggedright\arraybackslash}p{4.6cm}|*{10}{>{\centering\arraybackslash}p{0.55cm}|}@{}}
\caption{Project schedule by course week.}\label{tab:gantt}\\
\hline
\textbf{Task / milestone} & \textbf{2} & \textbf{3} & \textbf{4} &
\textbf{5} & \textbf{6} & \textbf{7} & \textbf{8} & \textbf{9} &
\textbf{10} & \textbf{11} \\
\hline
\endfirsthead
\hline
\textbf{Task / milestone} & \textbf{2} & \textbf{3} & \textbf{4} &
\textbf{5} & \textbf{6} & \textbf{7} & \textbf{8} & \textbf{9} &
\textbf{10} & \textbf{11} \\
\hline
\endhead
\hline
\endfoot
\hline
\endlastfoot
Scope and PA1 planning & $\blacksquare$ & $\blacksquare$ & & & & & & & & \\
PA1 submission & & $\blacklozenge$ & & & & & & & & \\
Requirements and acceptance criteria & & & $\blacksquare$ & $\blacksquare$ & & & & & & \\
PA2 submission & & & & $\blacklozenge$ & & & & & & \\
Architecture, database, and UI design & & & & & $\blacksquare$ & $\blacksquare$ & & & & \\
PA3 submission & & & & & & $\blacklozenge$ & & & & \\
PA4 test plan and environment (4.1--4.2) & & & & & & & $\blacksquare$ & & & \\
PA4 submission & & & & & & & $\blacklozenge$ & & & \\
Database and authentication (5.1.1) & & & & & & & $\blacksquare$ & & & \\
Booking backend preparation and implementation (5.1.2) & & & & & & & $\blacksquare$ & $\blacksquare$ & & \\
Room and booking-request UI (5.2.1--2) & & & & & & & $\blacksquare$ & $\blacksquare$ & & \\
Admin API (5.3.1) & & & & & & & $\blacksquare$ & $\blacksquare$ & & \\
Initial CI configuration (5.1.3) & & & & & & & $\blacksquare$ & & & \\
Admin screens (5.2.3) & & & & & & & & $\blacksquare$ & $\blacksquare$ & \\
Front end--back end integration (5.3.2) & & & & & & & & $\blacksquare$ & $\blacksquare$ & \\
Automation (5.4.2) & & & & & & & $\blacksquare$ & $\blacksquare$ & $\blacksquare$ & \\
Front end unit tests (5.2.4) & & & & & & & & $\blacksquare$ & $\blacksquare$ & \\
Back end unit tests and coverage completion (5.1.3) & & & & & & & & & $\blacksquare$ & \\
System testing and retesting of integrated workflows (5.4.3) & & & & & & & & $\blacksquare$ & $\blacksquare$ & \\
Test report and retrospective (5.4.4) & & & & & & & & $\blacksquare$ & $\blacksquare$ & \\
Defect fixing after system testing (5.3.3) & & & & & & & & & $\blacksquare$ & \\
Deployment and smoke test (5.1.4) & & & & & & & & & $\blacksquare$ & \\
W9 documentation draft; W10 final materials/video (5.3.4) & & & & & & & & $\blacksquare$ & $\blacksquare$ & \\
Continuous management (7.1.1--4) & $\blacksquare$ & $\blacksquare$ & $\blacksquare$ & $\blacksquare$ & $\blacksquare$ & $\blacksquare$ & $\blacksquare$ & $\blacksquare$ & $\blacksquare$ & \\
PA5 submission and final review & & & & & & & & & $\blacklozenge$ & \\
Final presentation / demonstration & & & & & & & & & & $\blacklozenge$ \\
\end{longtable}
\endgroup

\noindent\textbf{Legend:} Column numbers indicate course weeks; $\blacksquare$ = planned work; $\blacklozenge$ = milestone. The WBS allocation table gives exact weekly effort where a workstream spans multiple weeks.

\medskip
\noindent\textbf{Dependencies and sequencing:}
\begin{itemize}[leftmargin=*,itemsep=2pt]
  \item Finalize the database schema and API contracts from PA3 before implementation.
  \item Complete database and authentication work before integrating workflows that require sign-in.
  \item Integrate incrementally in W9--W10: connect student-booking and admin-approval flows when ready in W9, then connect room and maintenance management after the corresponding UI and API are complete in W10.
  \item Prepare test cases and automation in advance, but run workflow tests only on flows that have been integrated. In W10, complete the room/maintenance UI and API, integrate them, run the corresponding system tests, fix defects and retest, then verify the release.
  \item Finalize the test report and PA5 video using the deployed build and results from completed checks.
\end{itemize}
\textbf{Expected critical path.} Under the current week-level planning
assumptions, the expected sequence is requirements baseline (PA2)
$\rightarrow$ architecture, database design, and API contracts (PA3)
$\rightarrow$ database/authentication and core booking backend (5.1.1--5.1.2)
$\rightarrow$ front-end--back-end integration (5.3.2)
$\rightarrow$ system testing, defect fixing, and retesting (5.4.3, 5.3.3)
$\rightarrow$ deployment and smoke testing (5.1.4)
$\rightarrow$ final PA5 report, release documentation, and video (5.4.4, 5.3.4)
$\rightarrow$ final demonstration. Delayed required inputs can postpone
integration, testing, and handover. Integration is planned in parts across W9
and W10. Room/maintenance UI and API work must be complete and integrated before
full system testing of those functions; in W10, sequence their completion,
integration, testing, defect fixing/retesting, and release verification by these
prerequisites. The final report and video use the deployed build after checks
are complete. Work may overlap partially when its prerequisites are available.
This is an expected path, not a CPM or float calculation; confirm it when
dependencies and durations are detailed.

\subsection{Roles and Responsibilities}

The responsibilities below follow the team's agreed role allocation. All members
share responsibility for reviewing the report, testing the integrated application,
communicating progress, and recording their own contributions.

\begingroup
\small
\setlength{\tabcolsep}{5pt}
\renewcommand{\arraystretch}{1.22}
\begin{center}
\begin{tabular}{@{}>{\raggedright\arraybackslash}p{3.0cm}
>{\centering\arraybackslash}p{1.65cm}
>{\centering\arraybackslash}p{1.65cm}
>{\raggedright\arraybackslash}p{6.0cm}@{}}
\hline
\textbf{Role} & \textbf{Main} & \textbf{Backup} & \textbf{Responsibilities} \\
\hline
Team Lead / Project Manager & Long & Thong &
Maintain the schedule and task board, assign and follow up tasks, organise weekly
status meetings, track risks, and coordinate submissions and communication. \\
\hline
Business Analyst / Requirements & Thong & Long &
Gather and document requirements, booking rules, use cases, and acceptance criteria;
check consistency across PA documents and coordinate requirement reviews. \\
\hline
Backend Developer / Database / Architecture & Phuc & Ngan &
Lead C4 and architecture work, ADRs, and ERD; design and maintain the database and
backend; implement authentication, role checks, booking validation, and API endpoints. \\
\hline
Frontend Developer / UI Designer & Ngan & Phuc &
Design and implement Student and Administrator screens; display room availability and
booking statuses; connect the interface to backend APIs and check usability. \\
\hline
Testing and Quality Coordination & Thong & Long &
Lead test planning, test data, workflow and integration testing, defect tracking,
retesting, and test evidence/reporting; all members support testing as assigned. \\
\hline
\end{tabular}
\end{center}
\endgroup

\subsection{Human Resources and Cost Estimate}

The project is planned for a four-member student team. The baseline workload is
236 person-hours according to the WBS. For a comparable cost estimate, the team uses a notional rate of
VND 100,000 per person-hour. This rate is a planning assumption to represent the
value of effort, not an actual payment to team members. Hosting and contingency
figures are provisional and should be revised once the deployment approach is known.

\begingroup
\small
\setlength{\tabcolsep}{6pt}
\renewcommand{\arraystretch}{1.22}
\begin{center}
\begin{tabular}{@{}>{\raggedright\arraybackslash}p{8.5cm}
>{\centering\arraybackslash}p{3.4cm}@{}}
\hline
\textbf{Cost item} & \textbf{Estimated cost (VND)} \\
\hline
Human effort: 236 person-hours $\times$ VND 100,000/hour & 23,600,000 \\
\hline
Development tools and collaboration software (free/student tiers) & 0 \\
\hline
Hosting / deployment allowance for the course demonstration & 500,000 \\
\hline
Other expenses / contingency allowance & 500,000 \\
\hline
\textbf{Total estimated project cost} & \textbf{24,600,000} \\
\hline
\end{tabular}
\end{center}
\endgroup

The 500,000 VND contingency is a monetary reserve, separate from capacity headroom
and not an additional labor estimate. The estimate assumes that the team uses available free or student plans for source
control, documentation, and development. It excludes the purchase of personal
computers, existing internet access, and actual salary payments. The team will
review the estimate if the feature scope, hosting option, or course requirements change.
